% ---------------------------------------------------------------------------
% ---------------------------------------------------------------------------
% Modelo LaTex para preparação do documento final de Monografia TCC
% O modelo está em conformidade com ABNT NBR
% Faculdade do Piaui
% ---------------------------------------------------------------------------
% ---------------------------------------------------------------------------

\documentclass[
	% -- opções da classe memoir --
	12pt,					% tamanho da fonte
	openright,				% capítulos começam em pág ímpar (insere página vazia caso preciso)
	twoside,					% para impressão em verso e anverso. Oposto a oneside
	a4paper,					% tamanho do papel. 
	% -- opções da classe abntex2 --
	%chapter=TITLE,			% títulos de capítulos convertidos em letras maiúsculas
	%section=TITLE,			% títulos de seções convertidos em letras maiúsculas
	%subsection=TITLE,		% títulos de subseções convertidos em letras maiúsculas
	%subsubsection=TITLE,	% títulos de subsubseções convertidos em letras maiúsculas
	% -- opções do pacote babel --
	english,					% idioma adicional para hifenização
	%french,					% idioma adicional para hifenização
	%spanish,				% idioma adicional para hifenização
	brazil					% o último idioma é o principal do documento
	]{abntex2}

% ---------------------
% Pacotes OBRIGATÓRIOS
% ---------------------
\usepackage{lmodern}				% Usa a fonte Latin Modern			
\usepackage[T1]{fontenc}			% Selecao de codigos de fonte.
\usepackage[utf8]{inputenc}		% Codificacao do documento (conversão automática dos acentos)
\usepackage{lastpage}			% Usado pela Ficha catalográfica
\usepackage{indentfirst}			% Indenta o primeiro parágrafo de cada seção.
\usepackage{color}				% Controle das cores
\usepackage{graphicx,graphicx}	% Inclusão de gráficos
\usepackage{epsfig,subfig}		% Inclusão de figuras
\usepackage{microtype} 			% Melhorias de justificação
% ---------------------
		
% ---------------------
% Pacotes ADICIONAIS
% ---------------------
\usepackage{lipsum}						% Geração de dummy text
\usepackage{amsmath,amssymb,mathrsfs}	% Comandos matemáticos avançados 
\usepackage{setspace}  					% Para permitir espaçamento simples, 1 1/2 e duplo
\usepackage{verbatim}					% Para poder usar o ambiente "comment"
\usepackage{tabularx} 					% Para poder ter tabelas com colunas de largura auto-ajustável
\usepackage{afterpage} 					% Para executar um comando depois do fim da página corrente
\usepackage{url} 						% Para formatar URLs (endereços da Web)
% ---------------------

% ---------------------
% Pacotes de CITAÇÕES
% ---------------------
\usepackage[brazilian,hyperpageref]{backref}	% Paginas com as citações na bibl
\usepackage[alf]{abntex2cite}				% Citações padrão ABNT (alfa)
%\usepackage[num]{abntex2cite}				% Citações padrão ABNT (numericas)
% ---------------------

% Configurações de CITAÇÕES para abntex2
% --- 
% CONFIGURAÇÕES DE PACOTES
% --- 

% ---
% Configurações do pacote backref
% Usado sem a opção hyperpageref de backref
\renewcommand{\backrefpagesname}{Citado na(s) página(s):~}
% Texto padrão antes do número das páginas
\renewcommand{\backref}{}
% Define os textos da citação
\renewcommand*{\backrefalt}[4]{
	\ifcase #1 %
		Nenhuma citação no texto.%
	\or
		Citado na página #2.%
	\else
		Citado #1 vezes nas páginas #2.%
	\fi}%
% ---

% Inclusão de dados para CAPA e FOLHA DE ROSTO (título, autor, orientador, etc.)
% ---
% Informações de dados para CAPA e FOLHA DE ROSTO
% ---
\titulo{CONTROLE E MONITORAMENTO DE DEMANDA ENERGÉTICA}
\autor{João Marcos Oliveira Lima \\ Antonio Idalécio Pereira de Souza}
\local{Dourados - MS}
\data{Xxxx de 20XX}
\orientador{Ricardo Camparim}
\instituicao{%
  Universidade Federal da Grande Dourados - UFGD}
\tipotrabalho{Trabalho de Conclusão de Curso (TCC)}
% O preambulo deve conter o tipo do trabalho, o objetivo,
% o nome da instituição e a área de concentração
\preambulo{\textbf{Trabalho de Conclusão de Curso} Monografia apresentada a Universidade Federal da Grande Dourados (UFGD) como um dos pré-requisitos para a obtenção do grau de Bacharel em Ciências da Computação.}
% ---

% Inclui Configurações de aparência do PDF Final
%  Configurações de aparência do PDF final
% NÃO ALTERAR!!!

% alterando o aspecto da cor azul
\definecolor{blue}{RGB}{41,5,195}

% informações do PDF
\makeatletter
\hypersetup{
     	%pagebackref=true,
		pdftitle={\@title}, 
		pdfauthor={\@author},
    		pdfsubject={\imprimirpreambulo},
	    pdfcreator={LaTeX with abnTeX2},
		pdfkeywords={abnt}{latex}{abntex}{abntex2}{trabalho acadêmico}, 
		colorlinks=true,       		% false: boxed links; true: colored links
    		linkcolor=blue,          	% color of internal links
    		citecolor=blue,        		% color of links to bibliography
    		filecolor=magenta,      		% color of file links
		urlcolor=blue,
		bookmarksdepth=4
} 
\makeatother
% --- 

% O tamanho da identação do parágrafo é dado por:
\setlength{\parindent}{1.3cm}

% Controle do espaçamento entre um parágrafo e outro:
\setlength{\parskip}{0.2cm}  % tente também \onelineskip

% ---------------------
% Compila o indice
% ---------------------
\makeindex
% ---------------------

%%%%%%%%%%%%%%%%%%%%%%%%%%%
%%  INICIO DO DOCUMENTO  %%
%%%%%%%%%%%%%%%%%%%%%%%%%%%
\begin{document}

% Retira espaço extra obsoleto entre as frases.
\frenchspacing

% ----------------------------------------------------------
% ELEMENTOS PRÉ-TEXTUAIS (Capa, Resumo, Abstract, etc.)
% ----------------------------------------------------------
\pretextual

% Capa
% ---
% Impressão da Capa
% ---
  \begin{capa}%
    \begin{figure}[h!]%
        \centering%
        \includegraphics[scale=0.4]{figs/logo_ufgd.jpg}%
      \end{figure}%
    \center
	\ABNTEXchapterfont\large    {\large UFGD - UNIVERSIDADE FEDERAL DA GRANDE DOURADOS} \\
    
    {\large FACET - FACULDADE DE CIÊNCIAS EXATAS E TECNOLOGIA}  \\ 

    {\large BACHARELADO EM ENGENHARIA DE COMPUTAÇÃO}  \\ 
	%\vspace{1.5cm}

    \vfill
    \ABNTEXchapterfont\bfseries\LARGE\imprimirtitulo
    \vfill

	%\vfill
	\ABNTEXchapterfont\large\imprimirautor
	\vfill
%
    \large\imprimirlocal, \large\imprimirdata

    \vspace*{1cm}
  \end{capa}
% ---

% Folha de rosto (o * indica que haverá a ficha bibliográfica)
\imprimirfolhaderosto*

% Imprimir Ficha Catalografica
%% ---
% Ficha Catalográfica
% ---
% Isto é um exemplo de Ficha Catalográfica, ou ``Dados internacionais de
% catalogação-na-publicação''. Você pode utilizar este modelo como referência. 
% Porém, provavelmente a biblioteca da sua universidade lhe fornecerá um PDF
% com a ficha catalográfica definitiva após a defesa do trabalho. Quando estiver
% com o documento, salve-o como PDF no diretório do seu projeto e substitua todo
% o conteúdo de implementação deste arquivo pelo comando abaixo:
%
% \begin{fichacatalografica}
%     \includepdf{fig_ficha_catalografica.pdf}
% \end{fichacatalografica}
\begin{fichacatalografica}
	\vspace*{\fill}					% Posição vertical
	\hrule							% Linha horizontal
	\begin{center}					% Minipage Centralizado
	\begin{minipage}[c]{12.5cm}		% Largura
	
	\imprimirautor
	
	\hspace{0.5cm} \imprimirtitulo  / \imprimirautor. --
	\imprimirlocal, \imprimirdata-
	
	\hspace{0.5cm} \pageref{LastPage} p. : il. (algumas color.) ; 30 cm.\\
	
	\hspace{0.5cm} \imprimirorientadorRotulo~\imprimirorientador\\
	
	\hspace{0.5cm}
	\parbox[t]{\textwidth}{\imprimirtipotrabalho~--~\imprimirinstituicao,
	\imprimirdata.}\\
	
	\hspace{0.5cm}
		1. Palavra-chave1.
		2. Palavra-chave2.
		I. Orientador.
		II. Universidade xxx.
		III. Faculdade de xxx.
		IV. Título\\ 			
	
	\hspace{8.75cm} CDU 02:141:005.7\\
	
	\end{minipage}
	\end{center}
	\hrule
\end{fichacatalografica}
% ---

% Inserir Folha de Aprovação
% ---
% Assinaturas
% ---
% Isto é um exemplo de Folha de aprovação, elemento obrigatório da NBR
% 14724/2011 (seção 4.2.1.3). Você pode utilizar este modelo até a aprovação
% do trabalho. Após isso, substitua todo o conteúdo deste arquivo por uma
% imagem da página assinada pela banca com o comando abaixo:
%
% \includepdf{folhadeaprovacao_final.pdf}
%
\begin{folhadeaprovacao}

  \begin{center}
    {\ABNTEXchapterfont\large\imprimirautor}

    \vspace*{\fill}\vspace*{\fill}
    \begin{center}
      \ABNTEXchapterfont\bfseries\Large\imprimirtitulo
    \end{center}
    \vspace*{\fill}
    
    \hspace{.45\textwidth}
    \begin{minipage}{.5\textwidth}
        \imprimirpreambulo
    \end{minipage}%
    \vspace*{\fill}
   \end{center}
        
   {\centering Trabalho aprovado. \imprimirlocal, 01 de janeiro de 2014:}

   \assinatura{\textbf{\imprimirorientador} \\ Orientador} 
   \assinatura{\textbf{Professor} \\ Convidado 1}
   \assinatura{\textbf{Professor} \\ Convidado 2}
   \assinatura{\textbf{Professor} \\ Convidado 3}
      
   \begin{center}
    \vspace*{0.5cm}
    {\large\imprimirlocal}
    \par
    {\large\imprimirdata}
    \vspace*{1cm}
  \end{center}
  
\end{folhadeaprovacao}
% ---

% Dedicatória
%% ---
% Dedicatória
% ---
\begin{dedicatoria}
   \vspace*{\fill}
   \centering
   \noindent
   \textit{ Aos meus pais XXXXXXXX e YYYYYYY, \\ por sempre estarem comigo em todos os momentos.} \vspace*{\fill}
\end{dedicatoria}
% ---

% Agradecimentos
% ---
% Agradecimentos
% ---
\begin{agradecimentos}

Agradeço a Deus.

Agradeço aos meus pais, XXXXX e YYYYY, por ...

Aos meus irmãos, por.....

Agradeço ao meu orientador, XXXXXXXXX, por todos os conselhos, pela paciência e ajuda nesse período.

Aos meus amigos ...

Aos professores ...

À XXXXXX pelo apoio financeiro para realização deste trabalho de pesquisa.

\end{agradecimentos}
%% ---

% Epígrafe
% ---
% Epígrafe
% ---
\begin{epigrafe}
    \vspace*{\fill}
	\begin{flushright}
		\textit{``Não sei o que, \\
		          não sei o que,\\
                  não sei o que lá.''\\
		          (Stark Antônio, 2018)}
	\end{flushright}
\end{epigrafe}
% ---

% Resumo e Abstract
% ---
% RESUMOS
% ---

% RESUMO em português
\setlength{\absparsep}{18pt} % ajusta o espaçamento dos parágrafos do resumo
\begin{resumo}
 Segundo a ABNT, o resumo deve ressaltar o
 objetivo, o método, os resultados e as conclusões do documento. A ordem e a extensão
 destes itens dependem do tipo de resumo (informativo ou indicativo) e do
 tratamento que cada item recebe no documento original. O resumo deve ser
 precedido da referência do documento, com exceção do resumo inserido no
 próprio documento. (\ldots) As palavras-chave devem figurar logo abaixo do
 resumo, antecedidas da expressão Palavras-chave:, separadas entre si por
 ponto e finalizadas também por ponto.

 \textbf{Palavras-chaves}: latex. abntex. editoração de texto.
\end{resumo}

% ABSTRACT in english
\begin{resumo}[Abstract]
 \begin{otherlanguage*}{english}
   This is the english abstract.

   \vspace{\onelineskip}
 
   \noindent 
   \textbf{Keywords}: latex. abntex. text editoration.
 \end{otherlanguage*}
\end{resumo}

% Lista de ilustrações
\pdfbookmark[0]{\listfigurename}{lof}
\listoffigures*
\cleardoublepage

% Lista de tabelas
\pdfbookmark[0]{\listtablename}{lot}
\listoftables*
\cleardoublepage

% Lista de abreviaturas e siglas
\begin{siglas}
  \item[ABNT] Associação Brasileira de Normas Técnicas
  \item[abnTeX] ABsurdas Normas para TeX
\end{siglas}

% Lista de símbolos
\begin{simbolos}
  \item[$ \Gamma $] Letra grega Gama
  \item[$ \Lambda $] Lambda
  \item[$ \zeta $] Letra grega minúscula zeta
  \item[$ \in $] Pertence
\end{simbolos}

% Inserir o SUMÁRIO
\pdfbookmark[0]{\contentsname}{toc}
\tableofcontents*
\cleardoublepage

% ----------------------------------------------------------
% ELEMENTOS TEXTUAIS (Capítulos)
% ----------------------------------------------------------
\textual
% Elementos textuais com numeração arábica
\pagenumbering{arabic}
% Reinicia a contagem do número de páginas
\setcounter{page}{1}

% Inclui cada capitulo da Dissertação
% ----------------------------------------------------------
% Introdução 
% Capítulo sem numeração, mas presente no Sumário
% ----------------------------------------------------------

\chapter[Introdução]{Introdução}
%\addcontentsline{toc}{chapter}{Introdução}

Ao longo da história da humanidade, foram presenciadas várias formas de geração e/ou utilização de energia para diversos fins, desde a utilização da força braçal do homem e de animais, do aproveitamento das forças da água e do vento, da utilização do vapor, até a descoberta e utilização da energia elétrica. Todas estas formas de utilização de energia tiveram grande importância em suas respectivas épocas e foram de extrema importância para o desenvolvimento da humanidade, o mesmo ocorre na atualidade com a utilização da energia elétrica que se tornou uma das principais fontes de energia no mundo, a qual já proporcionou inúmeros avanços em diversas áreas como transporte, industrial, médicas, comunicação, entre outras \cite{jones1991electrical}.

Vários fatores contribuíram para o crescimento desta como fonte principal de energia, dentre os quais podemos destacar, a possibilidade de transformação desta em praticamente qualquer outra forma de energia e a relativa facilidade de seu transporte em grandes distâncias através de linhas de alta tensão. No entanto, no Brasil, o setor energético vem enfrentando uma crise de investimentos, e a escassez de chuvas nas regiões onde se encontram os maiores reservatórios de água essenciais para a geração de energia elétrica pelas usinas hidrelétricas, além do aumento da demanda por parte dos consumidores, consequentemente causa a necessidade da utilização das termoelétricas, esses fatores podem ser citados como contribuintes para o aumento gradativo no custo da produção de energia \cite[p. 11]{alissonUTFPR2016}. 

%A energia elétrica pode ser transformada de maneira direta em qualquer outra energia seja ela mecânica, química, térmica, entre outras. Além de apresentar facilidade de transporte e grande alcance através das linhas de alta tensão. Dessa forma a energia elétrica se tornou uma das fontes energéticas mais utilizadas no século XX, o que acabou provocando um crescimento vertiginoso na sua produção. No entanto no Brasil o setor energético vem enfrentando uma crise de investimentos, e a escassez de chuvas nas regiões onde se encontram os maiores reservatórios de água essenciais para a geração de energia elétrica pelas usinas hidrelétricas, além do aumento da demanda por parte dos consumidores, consequentemente causa a necessidade da utilização das termoelétricas, esses fatores podem ser citados como contribuintes para o aumento gradativo no custo da produção de energia \cite[p. 11]{alissonUTFPR2016}.

Esse aumento no custo da produção causado pela necessidade do uso da geração térmica, implica no aumento das tarifas de energia elétrica. Por mais que as usinas térmicas apresentem um custo de geração mais elevado, representam a segurança do abastecimento, e funcionam como uma fonte alternativa de suplementação do sistema quanto por motivos de escassez de chuvas as hidrelétricas não possuem condições suficientes para gerar toda a energia que o país necessita \cite{tancredi2013brasil}. Hoje é possível notar que a energia elétrica está presente em todos os setores da economia, então uma melhor gestão do consumo de energia torna-se uma tarefa interessante para fins de economia e redução de gastos.

% Em razão ao aumento no custo da produção de energia elétrica, o que consequentemente causa um aumento no preço repassado ao consumidor.

Normalmente os consumidores só possuem acesso limitado ao seu consumo de energia elétrica, isto é, os mesmos só possuem acesso a estas informações no final do mês. Além disso o usuário também não possui qualquer forma de realizar o monitoramento ou controle individual de cada um de seus equipamentos o que ofereceria uma melhor forma de controle sobre seus gastos, onde seria possível perceber os equipamentos com maior consumo energético e tomar alguma providência. Entretanto para tentar solucionar este problema o mercado disponibiliza alguns sistemas que executam o monitoramento em tempo real como o \textit{OWL CM160}\footnote{https://www.efimarket.pt/monitor-energetico-owl-cm-160} e \textit{Energy Detective Pro}\footnote{http://www.theenergydetective.com/prodocs} que foram utilizados como base neste estudo, ambos possuem conexão \textit{wireless} e atendem também o sistema trifásico. No entanto, são soluções que necessitam ser importadas além de não possuírem suporte de tradução e estão a um preço consideravelmente alto. Nosso produto a ser desenvolvido visa a redução de custos de implementação deste sistema em relação aos disponíveis no mercado para que sejam acessíveis ao consumidor, além de  permitir o monitoramento em tempo real de cada um dos seus equipamentos, proporcionar ao usuário uma conexão remota de qualquer local, notificar o consumidor quando sua demanda de energia estiver próxima de ser alcançada e fornecer uma fácil instalação.

\section{Motivação}\label{sec:motivação}
% \addcontentsline{toc}{section}{Motivação}
A motivação deste trabalho advém da necessidade de permitir que o usuário possa monitorar e controlar o consumo de seus equipamentos remotamente através de uma aplicação móvel para smartphones, a fim de conscientizá-lo e permitir que este possa entender melhor seus vícios de consumo.

\section{Objetivos}\label{sec:objetivos}
\subsection{Objetivos Gerais}\label{sec:objetivos_gerais}
Desenvolver um equipamento capaz de medir o consumo de energia elétrica de residências e estabelecimentos em geral, bem como transmitir essas informações através da internet e disponibilizá-las em uma aplicação para \textit{smartphones} para que o usuário possa acompanhar o consumo de energia elétrica de sua residência ou estabelecimento de qualquer lugar.

\subsection{Objetivos Específicos}\label{sec:objetivos_especificos}

\begin{itemize}
	\item Pesquisar os sistemas de monitoramento de consumo de energia elétrica disponíveis no mercado;
	\item Desenvolver um protótipo do produto para monitoramento de demandas energéticas, integrando a aquisição de dados e a determinação da potência, englobando \textit{hardware} e \textit{software};
    \item Desenvolver um aplicativo que permita ao usuário acompanhar e controlar o seu consumo energético;
    \item Realizar a verificação e validação do produto por meio de experimentos em campo, comparando os valores adquiridos com os obtidos por instrumentos confiáveis de medição.
\end{itemize}

\section{Organização}\label{sec:organização}

%\addcontentsline{toc}{section}{Objetivos}


% PARTE - Define a divisão do documento em partes (Não é obrigatório)
%\part{Preparação da pesquisa}
% ---
%\chapter{Uso de referências bibliográficas}
% ---

%-
%\subsection{Acentuação de referências bibliográficas}
%-
\chapter{Estado da Arte}\label{cap:estArte}

\section{Trabalhos Relacionados a Isto}\label{sec:primTrab}
%\addcontentsline{toc}{section}{Trabalhos Relacionados a Isto}

\chapter{Materiais e Métodos}\label{cap:ferramentas}

\section{Considerações Finais}


% PARTE
%\part{Proposta}
\chapter{Sistema Proposto}\label{cap:proposta}

Esse trabalho propõe um sistema de... 


\section{Primeira Parte do Sistema Proposto}

\section{Considerações Finais}


% PARTE
%\part{Parte Final}
\chapter{Resultados e Discussão}\label{cap:resultados}

\section{Base de Dados}

\section{Considerações Finais}

\chapter{Conclusões e Trabalhos Futuros}\label{cap:conclusao}
%\addcontentsline{toc}{chapter}{Conclusão e Trabalhos Futuros}

\section{Conclusões}

\section{Trabalhos Futuros}


% ----------------------------------------------------------
% ELEMENTOS PÓS-TEXTUAIS (Referências, Glossário, Apêndices)
% ----------------------------------------------------------
\postextual

% Referências bibliográficas
\bibliography{bibliografia}

% Glossário (Consulte o manual)
%\glossary

% Apêndices
% ----------------------------------------------------------
% Apêndices
% ----------------------------------------------------------

% ---
% Inicia os apêndices
% ---
\begin{apendicesenv}

% Imprime uma página indicando o início dos apêndices
\partapendices

% ----------------------------------------------------------
\chapter{Primeiro Apêncice}
% ----------------------------------------------------------

% ----------------------------------------------------------
\chapter{Perceba que o texto do título desse segundo apêndice é bem grande}
% ----------------------------------------------------------

\end{apendicesenv}
% ---

% Anexos
% ----------------------------------------------------------
% Apêndices
% ----------------------------------------------------------

% ---
% Inicia os anexos
% ---
\begin{anexosenv}

% Imprime uma página indicando o início dos anexos
\partanexos

% ---
\chapter{Nome do Primeiro Anexo}
% ---

% ---
\chapter{Nome de Outro Anexo}
% ---

\end{anexosenv}

% Índice remissivo (Consultar manual)
%\phantompart
%\printindex

\end{document}
